\documentclass[pdf,allen]{prosper}
\usepackage{graphicx,url}

\renewcommand{\P}{\ P}
%\newcommand{\jc}{\makebox[0.4pt][l]{0}1}
\newcommand{\jc}{\oslash}

\title{}
\author{Changepoint detection in network measurements}
\institution{Allen B. Downey}

\begin{document}
\maketitle

\begin{slide}{}

Fundamental problem: 

\begin{itemize}

\item Predict next value in time series.

\end{itemize}

Applications:

\begin{itemize}

\item Protocol parameters (timeouts).

\item Resource selection, scheduling.

\item User feedback.

\end{itemize}


\end{slide}

\begin{slide}{}

Two kinds of prediction:

\begin{itemize}

\item Single value prediction.

\item Predictive distribution.

    \begin{itemize}

    \item Summary stats.
    \item Intervals.
    \item $P(error > thresh)$
    \item $E[cost(error)]$

    \end{itemize}


\end{itemize}


\end{slide}

\begin{slide}{}

If we assume stationarity, life is good.

\begin{itemize}

\item Accumulate data indefinitely.

\item Predictive distribution = observed distribution.

\end{itemize}


\end{slide}

\begin{slide}{}

Non-stationary models:

\begin{itemize}

\item Trends + noise. 

\item Level + changepoint + noise.

\end{itemize}


\end{slide}

\begin{slide}{}

Network performance:

\begin{itemize}

\item Some trends (accumulating queue). 

\item Many abrupt changepoints.

\begin{itemize}

\item Beginning and end of transfers.

\item Routing changes.

\item Hardware failure, replacement.

\end{itemize}

\end{itemize}


\end{slide}

\begin{slide}{}

Prediction with {\em known} changepoints:

\begin{itemize}

\item Use data back to the latest changepoint.

\item Less accurate immediately after.

\end{itemize}


\end{slide}

\begin{slide}{}

Prediction with {\em probablistic} changepoints.


\[P(i) = \mbox{prob of a changepoint after point}\ i\]

Example:

\begin{itemize}

\item 150 data points.

\item $P(50) = 0.7$

\item $P(100) = 0.5$

\end{itemize}

How do we generate a predictive distribution?


\end{slide}

\begin{slide}{}

Two steps:

\begin{itemize}

\item Derive $P(i+)$ = prob that $i$ is the last changepoint

\item Compute weighted mix going back to each $i$.

\end{itemize}

Example:

\begin{itemize}

\item $P(50) = 0.7 \quad \quad P(100) = 0.5$

\item $P(50+) = 0.35 \quad P(100+) = 0.5$

\item Plus 0.15 chance of no changepoint.

\end{itemize}

\end{slide}

\begin{slide}{}

Predictive distribution = 

\begin{eqnarray*}
0.50  &\cdot & edf(100, 150) \oplus \\
0.35  &\cdot & edf(50, 150) \oplus \\
0.15 &\cdot & edf(0, 150) \\
\end{eqnarray*}



\end{slide}


\begin{slide}{}

So how do we generate the probabilities $P(i+)$?

\vspace{0.1in}

Three steps:

\begin{itemize}

\item Bayes' theorem.

\item Simple case: we know there is 1 changepoint.

\item General case: unknown \# of changepoints.

\end{itemize}


\end{slide}


\begin{slide}{}

Bayes' theorem (diachronic interpretation)

\[ P(H|E) = \frac{P(E|H)}{P(E)} P(H) \]

\begin{itemize}

\item $H$ is a hypothesis, $E$ is a body of evidence.

\item $P(H|E)$: posterior

\item $P(H)$: prior

\item $P(E|H)$ is usually easy to compute.

\item $P(E)$ is often not.

\end{itemize}


\end{slide}


\begin{slide}{}

Unless we have a suite of exclusive hypotheses.

\[ P(E) = \sum_{H_i \in S} P(E|H_i) P(H_i) \]

In that case life is good.

\end{slide}


\begin{slide}{}

\begin{itemize}

\item If we know there there is exactly one
changepoint in an interval...

\item ...then the $P(i)$ are exclusive hypotheses,

\item and all we need is $P(E|i)$.

\end{itemize}

Which is pretty much a solved problem.

\end{slide}


\begin{slide}{}

What if the \# of changepoints is unknown?

\begin{itemize}

\item $P(i)$ are no longer exclusive.

\item But the $P(i+)$ are.

\item And we can write a system of equations for $P(i+)$.

\end{itemize}

\end{slide}


\begin{slide}{}

\[ P(i^+) = P(i^+ | \jc) \P(\jc) + \sum_{j<i} P(i^+ | j^{++}) \P(j^{++})  \]


\begin{itemize}

\item $P(j^{++})$ is the prob that the second-to last changepoint is at $i$.

\item $P(i^+ | j^{++})$ reduces to the simple problem.

\item $P(\jc)$ is the prob that we have not seen two changepoints.

\item $P(i^+ | \jc)$ reduces to the simple problem (plus).

\end{itemize}


Great, so what's $P(j^{++})$?

\end{slide}



\begin{slide}{}

\[ P(i^{++}) = \sum_{k>i} P(i^{++} | k^+) \P(k^+) \]

\begin{itemize}

\item $P(i^{++} | k^+)$ is just $P(i^+)$ computed at time $k$.

\item So we can solve for $P(i^+)$ in terms of $P(i^{++})$.

\item And $P(i^{++})$ in terms of $P(i^+)$.

\item Calling Dr. Jacobi!

\end{itemize}



\end{slide}



\begin{slide}{}

Implementation:

\begin{itemize}

\item Need to keep $n^2/2$ previous values.

\item And $n^2/2$ summary statistics.

\item And it takes $n^2$ work to do an update.

\item But, we only have to go back two changepoints,

\item ...so we can keep $n$ small.

\end{itemize}


\end{slide}


\begin{slide}{}

\raisebox{-2.8in}{
\includegraphics[width=3.0in,height=3.1in]{../../python/change/example.eps}
}
\nolinebreak
\begin{minipage}[t]{1.8in}

\begin{itemize}

\item Synthetic series with two changepoints.

\item $\mu = -0.5, 0.5, 0.0$

\item $\sigma = 1.0$

\item $P(\jc) = 0.04$

\end{itemize}

\end{minipage}
\end{slide}



\begin{slide}{}

\raisebox{-2.8in}{
\includegraphics[width=3.0in,height=3.1in]{../../python/change/nile.eps}
}
\nolinebreak
\begin{minipage}[t]{1.8in}

\begin{itemize}

\item The ubiquitous Nile dataset.

\item Change in 1898.

\item Estimated probs can be mercurial.

\end{itemize}

\end{minipage}
\end{slide}


\begin{slide}{}

\raisebox{-2.8in}{
\includegraphics[width=3.0in,height=3.1in]{../../python/change/var.eps}
}
\nolinebreak
\begin{minipage}[t]{1.8in}

\begin{itemize}

\item Can also detect change in variance.

\item $\mu = 1, 0, 0$

\item $\sigma = 1, 1, 0.5$

\item Estimated $P(i^+)$ is good.

\item Estimated $P(i^{++})$ less certain.

\end{itemize}

\end{minipage}
\end{slide}



\begin{slide}{}

\begin{itemize}

\item Qualitative behavior seems good. 

\item Quantitative tests:

\begin{itemize}

\item Compare to GLR for online alarm problem.

\item Test predictive distribution with synthetic data,

\item ... and with real data.

\end{itemize}

\end{itemize}

\end{slide}



\begin{slide}{}

Online alarm problem:

\begin{itemize}

\item Observe process in real time.

\item $\mu_0$ and $\sigma$ known.

\item $\tau$ and $\mu_1$ unknown.

\item Raise alarm when $f(data) > thresh$.

\item Minimize delay.

\item Minimize false alarm rate.

\end{itemize}


\end{slide}



\begin{slide}{}

GLR = generalized likelihood ratio.

\begin{itemize}

\item Compute decision function $g_k$.

\item $E[g_k] = 0$ before the changepoint,

\item ... increases after.

\item Alarm when $g_k > h$.

\item GLR is optimal when $\mu_1$ is known.

\end{itemize}


\end{slide}



\begin{slide}{}

CPP = change point probability

\[ P(changepoint) = \sum_{i=0}^{k} P(i^+) \]

\begin{itemize}

\item Alarm when $P(changepoint) > thresh$.

\end{itemize}

\end{slide}



\begin{slide}{}

\raisebox{-2.8in}{
\includegraphics[width=3.0in,height=3.1in]{../../python/change/matfar.eps}
}
\nolinebreak
\begin{minipage}[t]{1.8in}

\begin{itemize}

\item $\mu = 0, 1$

\item $\sigma = 1$

\item $\tau \sim \mbox{Exp}(0.01)$ 

\item Goodness = lower mean delay for same false alarm rate.

\end{itemize}

\end{minipage}
\end{slide}




\begin{slide}{}

\raisebox{-2.8in}{
\includegraphics[width=3.0in,height=3.1in]{../../python/change/matsig.eps}
}
\nolinebreak
\begin{minipage}[t]{1.8in}

\begin{itemize}

\item Fix false alarm rate = 5\%

\item Vary $\sigma$.

\item CPP does well with small $S/N$.

\end{itemize}

\end{minipage}
\end{slide}


\begin{slide}{}

So it works on a simple problem.

\vspace{0.1in}

Future work:

\begin{itemize}

\item Other changepoint problems (location, tracking, prediction).

\item Other data distributions (lognormal).

\item Testing robustness (real data).

\end{itemize}

\end{slide}


\begin{slide}{}

Good news:

\begin{itemize}

\item Very general framework.

\item Seems to work.

\item Many possible applications.

\end{itemize}

\end{slide}


\begin{slide}{}

Bad news:

\begin{itemize}

\item Still some wrinkles to iron.

\item $n^2$ space and time may be fatal.

\item May be overkill for original application.


\end{itemize}


\end{slide}


\begin{slide}{}

\begin{itemize}

\item More at {\tt http://allendowney.com/changepoint}

\item Or email {\tt downey@allendowney.com}


\end{itemize}


\end{slide}


\end{document}
